% GENERATED FILE. Edit canonical lessons, then run npm run generate:books.
% canonical-source-hash: fnv1a64:8beb7cdcbf214a12
% canonical-lessons: GE-C458-eichhoernchen, GE-C458-igel, GE-C458-wildschwein, GE-C458-hirsch, GE-C458-zwitschern

\chapter{Tiere im Wald}
\label{ch:ge-tiere-im-wald}
\hlchaptermodality{458}

\begin{chapteropening}
\textbf{By the end of this chapter:} \emph{I can name animals of the woods and their sounds: the squirrel, the hedgehog, the wild boar, the stag and birds chirping.}

Complete the last lesson of chapter 458: I can name animals of the woods and their sounds: the squirrel, the hedgehog, the wild boar, the stag and birds chirping.
\end{chapteropening}

\section[das Eichhörnchen]{das Eichhörnchen — the squirrel}

\label{lesson:GE-C458-eichhoernchen}



\begin{quote}
Before the new one: say the German for the bay, then the German for the shell.
\end{quote}

\subsection*{You'll want to know: das Eichhörnchen}
\textbf{das Eichhörnchen} — \textquotedblleft{}the squirrel\textquotedblright{}.

\begin{quote}
Ein Eichhörnchen sammelt Nüsse für den Winter.
\end{quote}

\begin{quote}
\emph{A squirrel is collecting nuts for the winter.}
\end{quote}

Same form in the plural. In Austria and the south also \emph{das Eichkätzchen}.

\begin{cousinweb}
\emph{Eichhörnchen} goes back to Old High German \emph{eihhorno} and Proto-Germanic \emph{aikwernô} 'squirrel' (EtymDB); later speakers heard \emph{Eiche} 'oak' and \emph{Horn} in it and added \emph{-chen} 'little'.
\end{cousinweb}

\begin{grammarlens}[title={What you've built}]
One more word for telling, explaining and describing at B1.
\end{grammarlens}

\subsection*{Your turn}
\begin{itemize}
\raggedright
  \item \emph{Say it:} \emph{das Eichhörnchen}
  \item \emph{Say it:} \emph{das Eichhörnchen}, once more
  \item \emph{Say it:} a sentence of your own with \emph{das Eichhörnchen}
  \item \emph{Recall:} say the German for the bay, then the German for the shell, then say \emph{das Eichhörnchen} again
\end{itemize}

\begin{tcolorbox}[breakable,colback=teal!4,colframe=teal!35!black,title={Before you move on}]
What does das Eichhörnchen mean? (\textquotedblleft{}The squirrel\textquotedblright{}.) Say it once more.
\end{tcolorbox}
\section[der Igel]{der Igel — the hedgehog}

\label{lesson:GE-C458-igel}



\begin{quote}
Before the new one: say the German for the shell, then the German for the squirrel.
\end{quote}

\subsection*{You'll want to know: der Igel}
\textbf{der Igel} — \textquotedblleft{}the hedgehog\textquotedblright{}.

\begin{quote}
Im Garten wohnt ein Igel unter dem Laub.
\end{quote}

\begin{quote}
\emph{A hedgehog lives under the leaves in the garden.}
\end{quote}

Same form in the plural.

\begin{cousinweb}
\emph{Igel} is inherited through Middle High German \emph{igel} from Old High German \emph{igil} and Proto-Germanic \emph{igilaz} (EtymDB); EtymDB lists Lithuanian \emph{ežỹs} 'hedgehog' as a distant cousin.
\end{cousinweb}

\begin{grammarlens}[title={What you've built}]
One more word for telling, explaining and describing at B1.
\end{grammarlens}

\subsection*{Your turn}
\begin{itemize}
\raggedright
  \item \emph{Say it:} \emph{der Igel}
  \item \emph{Say it:} \emph{der Igel}, once more
  \item \emph{Say it:} a sentence of your own with \emph{der Igel}
  \item \emph{Recall:} say the German for the shell, then the German for the squirrel, then say \emph{der Igel} again
\end{itemize}

\begin{tcolorbox}[breakable,colback=teal!4,colframe=teal!35!black,title={Before you move on}]
What does der Igel mean? (\textquotedblleft{}The hedgehog\textquotedblright{}.) Say it once more.
\end{tcolorbox}
\section[das Wildschwein]{das Wildschwein — the wild boar}

\label{lesson:GE-C458-wildschwein}



\begin{quote}
Before the new one: say the German for the squirrel, then the German for the hedgehog.
\end{quote}

\subsection*{You'll want to know: das Wildschwein}
\textbf{das Wildschwein} — \textquotedblleft{}the wild boar\textquotedblright{}.

\begin{quote}
Nachts kommen die Wildschweine bis an das Dorf.
\end{quote}

\begin{quote}
\emph{At night the wild boar come right up to the village.}
\end{quote}

Plural: \emph{die Wildschweine}.

\begin{cousinweb}
\emph{Wildschwein} is built from \emph{wild} 'wild' and \emph{Schwein} 'pig', both inherited: \emph{Schwein} goes back to Proto-Germanic \emph{swīną} and is a cousin of English \emph{swine} (EtymDB).
\end{cousinweb}

\begin{grammarlens}[title={What you've built}]
One more word for telling, explaining and describing at B1.
\end{grammarlens}

\subsection*{Your turn}
\begin{itemize}
\raggedright
  \item \emph{Say it:} \emph{das Wildschwein}
  \item \emph{Say it:} \emph{das Wildschwein}, once more
  \item \emph{Say it:} a sentence of your own with \emph{das Wildschwein}
  \item \emph{Recall:} say the German for the squirrel, then the German for the hedgehog, then say \emph{das Wildschwein} again
\end{itemize}

\begin{tcolorbox}[breakable,colback=teal!4,colframe=teal!35!black,title={Before you move on}]
What does das Wildschwein mean? (\textquotedblleft{}The wild boar\textquotedblright{}.) Say it once more.
\end{tcolorbox}
\section[der Hirsch]{der Hirsch — the stag, the deer}

\label{lesson:GE-C458-hirsch}



\begin{quote}
Before the new one: say the German for the hedgehog, then the German for the wild boar.
\end{quote}

\subsection*{You'll want to know: der Hirsch}
\textbf{der Hirsch} — \textquotedblleft{}the stag, the deer\textquotedblright{}.

\begin{quote}
Im Herbst hört man die Hirsche im Wald rufen.
\end{quote}

\begin{quote}
\emph{In autumn you can hear the stags calling in the forest.}
\end{quote}

Plural: \emph{die Hirsche}. The female is \emph{die Hirschkuh}.

\begin{cousinweb}
\emph{Hirsch} is inherited through Middle High German \emph{hirz} from Old High German \emph{hiruz} and Proto-Germanic \emph{herutaz} (EtymDB); English \emph{hart} 'male deer' is a cousin.
\end{cousinweb}

\begin{grammarlens}[title={What you've built}]
One more word for telling, explaining and describing at B1.
\end{grammarlens}

\subsection*{Your turn}
\begin{itemize}
\raggedright
  \item \emph{Say it:} \emph{der Hirsch}
  \item \emph{Say it:} \emph{der Hirsch}, once more
  \item \emph{Say it:} a sentence of your own with \emph{der Hirsch}
  \item \emph{Recall:} say the German for the hedgehog, then the German for the wild boar, then say \emph{der Hirsch} again
\end{itemize}

\begin{tcolorbox}[breakable,colback=teal!4,colframe=teal!35!black,title={Before you move on}]
What does der Hirsch mean? (\textquotedblleft{}The stag, the deer\textquotedblright{}.) Say it once more.
\end{tcolorbox}
\section[zwitschern]{zwitschern — to chirp, to twitter}

\label{lesson:GE-C458-zwitschern}



\begin{quote}
Before the new one: say the German for the wild boar, then the German for the stag.
\end{quote}

\subsection*{You'll want to know: zwitschern}
\textbf{zwitschern} — \textquotedblleft{}to chirp, to twitter\textquotedblright{}.

\begin{quote}
Am Morgen zwitschern die Vögel vor dem Fenster.
\end{quote}

\begin{quote}
\emph{In the morning the birds chirp outside the window.}
\end{quote}

\emph{Das Zwitschern} — birdsong.

\begin{cousinweb}
\emph{zwitschern} is a sound word that imitates the calls of small birds; EtymDB lists the sense 'to twitter', and English \emph{twitter} is formed the same way.
\end{cousinweb}

\begin{grammarlens}[title={What you've built}]
One more word for telling, explaining and describing at B1.
\end{grammarlens}

\subsection*{Your turn}
\begin{itemize}
\raggedright
  \item \emph{Say it:} \emph{zwitschern}
  \item \emph{Say it:} \emph{zwitschern}, once more
  \item \emph{Say it:} a sentence of your own with \emph{zwitschern}
  \item \emph{Recall:} say the German for the wild boar, then the German for the stag, then say \emph{zwitschern} again
\end{itemize}

\begin{tcolorbox}[breakable,colback=teal!4,colframe=teal!35!black,title={Before you move on}]
What does zwitschern mean? (\textquotedblleft{}To chirp, to twitter\textquotedblright{}.) Say it once more.
\end{tcolorbox}
